\chapter{Coronary Bypass}
\index{Coronary Bypass}

\section*{Definition and Overview}
Coronary bypass surgery, formally called coronary artery bypass grafting (CABG), is an operation that creates detours around narrowed or blocked coronary arteries so that blood can once again reach the heart muscle. Healthy vessels from the chest, leg, or arm are used to construct these bypasses. The surgery is a major procedure performed under general anaesthetic, usually with the heart stopped and supported by a heart-lung machine. It is recommended when coronary disease is extensive, when the left main artery is narrowed, or when stenting is less suitable, and it reliably relieves angina and improves long-term survival in appropriate patients.

\section*{Epidemiology}
Coronary bypass is one of the most common major operations in many countries and worldwide, with thousands performed each year. Candidates are typically older adults with multi-vessel coronary disease, diabetes, or left main disease. The proportion of patients treated with surgery has declined relative to stenting for single-vessel disease, but surgery remains the preferred revascularisation strategy for complex and diffuse disease, particularly in people with diabetes, where trials have shown better long-term outcomes than stenting.

\section*{Aetiology and Pathophysiology}
The operation addresses the underlying problem of coronary atherosclerosis.

\begin{itemize}
    \item \textbf{Multi-vessel atherosclerosis:} significant narrowings in two or three coronary arteries
    \item \textbf{Left main disease:} narrowing of the main artery supplying the left side of the heart
    \item \textbf{Diffuse and calcified disease:} long or heavily calcified plaques that are poorly suited to stenting
    \item \textbf{Diabetes:} associated with more extensive, distal disease that favours surgery
    \item \textbf{Failed or unsuitable stenting:} recurrent narrowing or complex anatomy
    \item \textbf{Ischaemic cardiomyopathy:} reduced pump function from chronic ischaemia that may improve with revascularisation
\end{itemize}

\section*{Clinical Features}
Patients considered for bypass surgery typically have significant symptoms or high-risk anatomy.

\begin{itemize}
    \item Angina with exertion, or unstable angina at rest
    \item Shortness of breath and fatigue
    \item A recent or prior heart attack
    \item Reduced exercise capacity
    \item Abnormal stress testing
    \item Angiography showing severe or extensive coronary disease
    \item Symptoms that persist despite optimal medical therapy
\end{itemize}

\section*{Differential Diagnosis}
The decision to operate is made jointly by cardiologists and cardiac surgeons, weighing surgery against alternatives.

\begin{itemize}
    \item Percutaneous coronary intervention (stenting), preferred for limited disease
    \item Medical therapy alone for mild disease or poor surgical candidates
    \item Valve disease, which may require repair or replacement instead
    \item Heart failure from cardiomyopathy, which may not be improved by bypass
    \item Non-cardiac causes of chest pain
\end{itemize}

\section*{When to Seek Medical Care}
Patients being assessed for bypass surgery receive detailed counselling about the risks and benefits. Any new or worsening cardiac symptoms require prompt attention.

\textbf{Contact your local emergency services for:}
\begin{itemize}
    \item Chest pain at rest or lasting more than a few minutes
    \item Pain spreading to the arm, jaw, neck, or back
    \item Shortness of breath, sweating, nausea, or dizziness with chest pain
    \item New chest pain, fever, or a wound problem after surgery
\end{itemize}

\section*{Diagnosis and Investigations}
Thorough assessment is carried out before the operation.

\begin{itemize}
    \item \textbf{Coronary angiography:} defines the number, site, and severity of blockages
    \item \textbf{Echocardiography:} assesses pump function and valve disease
    \item \textbf{ECG:} detects prior heart attacks and rhythm problems
    \item \textbf{Blood tests:} kidney and liver function, full blood count, and coagulation
    \item \textbf{Chest X-ray and CT:} assessment of the lungs, aorta, and chest anatomy
    \item \textbf{Vascular assessment:} checks of the leg arteries for suitability of vein grafts
    \item \textbf{Fitness optimisation:} blood pressure, blood sugar, weight, and dental health review
\end{itemize}

\section*{Management}
The procedure is delivered by a cardiac surgical team with structured aftercare.

\begin{itemize}
    \item \textbf{Pre-operative clinic:} education, medicines optimisation, and smoking cessation
    \item \textbf{Operation:} performed under general anaesthetic; grafts are attached beyond each blockage
    \item \textbf{Heart-lung machine:} used when the heart is stopped for the grafting
    \item \textbf{Beating-heart technique:} used in selected patients to avoid the bypass machine
    \item \textbf{Intensive care:} the first night or two after surgery
    \item \textbf{Ward recovery:} mobilisation, physiotherapy, and pain management
    \item \textbf{Discharge planning:} medicines, wound care, and activity advice
    \item \textbf{Cardiac rehabilitation:} exercise and education for six to twelve weeks
\end{itemize}

\section*{Complications}
\begin{itemize}
    \item Bleeding or the need for re-operation
    \item Wound infection at the chest or graft sites
    \item Atrial fibrillation and other arrhythmias
    \item Stroke or temporary neurological disturbance
    \item Kidney injury
    \item Heart attack or low cardiac output in the early post-operative period
    \item Chest infection, fluid collections, or lung collapse
    \item Long-term graft narrowing and recurrent angina
    \item Post-operative confusion or memory changes, usually temporary
\end{itemize}

\section*{Prognosis and Outcomes}
Bypass surgery produces excellent symptom relief and, in high-risk patients, better survival than stenting or medical therapy alone. Operative mortality in elective fit patients is low, typically around 1--2\%. The internal chest artery graft remains patent in the vast majority of patients at ten years, while vein grafts have a higher long-term failure rate, reinforcing the importance of risk factor control. Most patients are back to normal activities within two to three months, and many remain free of angina for a decade or more.

\section*{Prevention and Self-Care}
\begin{itemize}
    \item Stop smoking permanently
    \item Take antiplatelet medicines, statins, and other prescribed medicines for life
    \item Attend cardiac rehabilitation and exercise regularly
    \item Eat a heart-healthy diet and keep weight, cholesterol, and blood sugar controlled
    \item Protect the healing breastbone: no heavy lifting for about six weeks
    \item Care for wounds as instructed and report any redness, discharge, or fever
    \item Keep follow-up appointments with the cardiologist and GP
    \item Recognise and act promptly on any new cardiac symptoms
\end{itemize}

\section*{Sources and Further Reading}
The following reputable sources provide further information for healthcare students and patients seeking reliable, current advice about coronary bypass surgery.

\begin{itemize}
    \item Heart Foundation Australia. \textit{Coronary artery bypass graft surgery.} https://www.heartfoundation.org.au/
    \item NHS. \textit{Coronary artery bypass graft (CABG).} https://www.nhs.uk/conditions/coronary-artery-bypass-graft-cabg/
    \item Mayo Clinic. \textit{Coronary bypass surgery.} https://www.mayoclinic.org/tests-procedures/coronary-bypass-surgery/
    \item MedlinePlus. \textit{Coronary artery bypass surgery.} https://medlineplus.gov/ency/article/002946.htm
    \item American Heart Association. \textit{Bypass surgery.} https://www.heart.org/
    \item MSD Manual. \textit{Coronary artery bypass grafting.} https://www.msdmanuals.com/
\end{itemize}
